\documentclass[11pt]{article}

% Shared notes settings for Elements of AIML lectures (UPES).
% Keep this file free of \documentclass to allow reuse via \input.

\usepackage[utf8]{inputenc}
\usepackage[T1]{fontenc}
\usepackage[a4paper,margin=1in]{geometry}
\usepackage{hyperref}
\usepackage{xurl}
\usepackage{booktabs}
\usepackage{amsmath}
\usepackage{amssymb}
\usepackage{listings}
\usepackage{xcolor}
\usepackage{enumitem}
\usepackage{parskip}

% Code listings (general-purpose).
\lstset{
  basicstyle=\ttfamily\small,
  keywordstyle=\color{blue},
  commentstyle=\color{gray},
  stringstyle=\color{teal},
  showstringspaces=false,
  columns=fullflexible,
  frame=single,
  framerule=0.2pt,
  breaklines=true
}

\setlist[itemize]{topsep=0.3em,itemsep=0.2em}
\setlist[enumerate]{topsep=0.3em,itemsep=0.2em}

% Simple per-section source line.
\newcommand{\notesources}[1]{\par\noindent{\small\color{gray}\textit{Sources:} \sloppy #1\par}}

% Unit-wise source strings for this subject (used by slides/notes).
% Path is relative to the lecture `latex/` folder (the compilation CWD).
\IfFileExists{../../../_shared/aiml-sources.tex}{%
  % Source strings for Elements of AIML (used by slides + notes).

\newcommand{\AIMLRepoURL}{https://github.com/tali7c/Elements-of-AIML}

\newcommand{\AIMLLabSources}{%
Provost and Fawcett, \textit{Data Science for Business}.;%
McKinney, \textit{Python for Data Analysis}.;%
WEKA Documentation (Waikato Environment for Knowledge Analysis), accessed 2026-02-28.;%
Matplotlib and Seaborn documentation, accessed 2026-02-28.%
}

\newcommand{\AIMLUnitISources}{%
Rich and Knight, \textit{Artificial Intelligence}, McGraw Hill, 2017.;%
Russell and Norvig, \textit{Artificial Intelligence: A Modern Approach} (4e), Ch. 1--2 (Introduction, Intelligent Agents).;%
Poole and Mackworth, \textit{Artificial Intelligence: Foundations of Computational Agents} (2e), selected sections.%
}

\newcommand{\AIMLUnitIISources}{%
Russell and Norvig, \textit{Artificial Intelligence: A Modern Approach} (4e), Ch. 7--9 (Logical Agents, First-Order Logic, Inference).;%
Huth and Ryan, \textit{Logic in Computer Science} (2e), selected sections on propositional and first-order logic.;%
Genesereth and Nilsson, \textit{Logical Foundations of Artificial Intelligence}, selected chapters.%
}

\newcommand{\AIMLUnitIIISources}{%
Mueller and Massaron, \textit{Machine Learning for Dummies}, For Dummies, 2016.;%
Mitchell, \textit{Machine Learning}, selected chapters on learning basics and evaluation.;%
Scikit-learn documentation, accessed 2026-02-28.%
}

\newcommand{\AIMLUnitIVSources}{%
Mueller and Massaron, \textit{Machine Learning for Dummies}, For Dummies, 2016.;%
James, Witten, Hastie, and Tibshirani, \textit{An Introduction to Statistical Learning} (2e), selected chapters.;%
Scikit-learn documentation, accessed 2026-02-28.%
}

\newcommand{\AIMLUnitVSources}{%
NIST, \textit{AI Risk Management Framework (AI RMF 1.0)}, 2023.;%
Barocas, Hardt, and Narayanan, \textit{Fairness and Machine Learning} (online book), selected sections.;%
Knaflic, \textit{Storytelling with Data}, selected chapters on communicating results.%
}
%
}{%
  % Faculty copies live under `private/` so their compile CWD is different.
  \IfFileExists{../../../../public/_shared/aiml-sources.tex}{%
    % Source strings for Elements of AIML (used by slides + notes).

\newcommand{\AIMLRepoURL}{https://github.com/tali7c/Elements-of-AIML}

\newcommand{\AIMLLabSources}{%
Provost and Fawcett, \textit{Data Science for Business}.;%
McKinney, \textit{Python for Data Analysis}.;%
WEKA Documentation (Waikato Environment for Knowledge Analysis), accessed 2026-02-28.;%
Matplotlib and Seaborn documentation, accessed 2026-02-28.%
}

\newcommand{\AIMLUnitISources}{%
Rich and Knight, \textit{Artificial Intelligence}, McGraw Hill, 2017.;%
Russell and Norvig, \textit{Artificial Intelligence: A Modern Approach} (4e), Ch. 1--2 (Introduction, Intelligent Agents).;%
Poole and Mackworth, \textit{Artificial Intelligence: Foundations of Computational Agents} (2e), selected sections.%
}

\newcommand{\AIMLUnitIISources}{%
Russell and Norvig, \textit{Artificial Intelligence: A Modern Approach} (4e), Ch. 7--9 (Logical Agents, First-Order Logic, Inference).;%
Huth and Ryan, \textit{Logic in Computer Science} (2e), selected sections on propositional and first-order logic.;%
Genesereth and Nilsson, \textit{Logical Foundations of Artificial Intelligence}, selected chapters.%
}

\newcommand{\AIMLUnitIIISources}{%
Mueller and Massaron, \textit{Machine Learning for Dummies}, For Dummies, 2016.;%
Mitchell, \textit{Machine Learning}, selected chapters on learning basics and evaluation.;%
Scikit-learn documentation, accessed 2026-02-28.%
}

\newcommand{\AIMLUnitIVSources}{%
Mueller and Massaron, \textit{Machine Learning for Dummies}, For Dummies, 2016.;%
James, Witten, Hastie, and Tibshirani, \textit{An Introduction to Statistical Learning} (2e), selected chapters.;%
Scikit-learn documentation, accessed 2026-02-28.%
}

\newcommand{\AIMLUnitVSources}{%
NIST, \textit{AI Risk Management Framework (AI RMF 1.0)}, 2023.;%
Barocas, Hardt, and Narayanan, \textit{Fairness and Machine Learning} (online book), selected sections.;%
Knaflic, \textit{Storytelling with Data}, selected chapters on communicating results.%
}
%
  }{%
    % Source strings for Elements of AIML (used by slides + notes).

\newcommand{\AIMLRepoURL}{https://github.com/tali7c/Elements-of-AIML}

\newcommand{\AIMLLabSources}{%
Provost and Fawcett, \textit{Data Science for Business}.;%
McKinney, \textit{Python for Data Analysis}.;%
WEKA Documentation (Waikato Environment for Knowledge Analysis), accessed 2026-02-28.;%
Matplotlib and Seaborn documentation, accessed 2026-02-28.%
}

\newcommand{\AIMLUnitISources}{%
Rich and Knight, \textit{Artificial Intelligence}, McGraw Hill, 2017.;%
Russell and Norvig, \textit{Artificial Intelligence: A Modern Approach} (4e), Ch. 1--2 (Introduction, Intelligent Agents).;%
Poole and Mackworth, \textit{Artificial Intelligence: Foundations of Computational Agents} (2e), selected sections.%
}

\newcommand{\AIMLUnitIISources}{%
Russell and Norvig, \textit{Artificial Intelligence: A Modern Approach} (4e), Ch. 7--9 (Logical Agents, First-Order Logic, Inference).;%
Huth and Ryan, \textit{Logic in Computer Science} (2e), selected sections on propositional and first-order logic.;%
Genesereth and Nilsson, \textit{Logical Foundations of Artificial Intelligence}, selected chapters.%
}

\newcommand{\AIMLUnitIIISources}{%
Mueller and Massaron, \textit{Machine Learning for Dummies}, For Dummies, 2016.;%
Mitchell, \textit{Machine Learning}, selected chapters on learning basics and evaluation.;%
Scikit-learn documentation, accessed 2026-02-28.%
}

\newcommand{\AIMLUnitIVSources}{%
Mueller and Massaron, \textit{Machine Learning for Dummies}, For Dummies, 2016.;%
James, Witten, Hastie, and Tibshirani, \textit{An Introduction to Statistical Learning} (2e), selected chapters.;%
Scikit-learn documentation, accessed 2026-02-28.%
}

\newcommand{\AIMLUnitVSources}{%
NIST, \textit{AI Risk Management Framework (AI RMF 1.0)}, 2023.;%
Barocas, Hardt, and Narayanan, \textit{Fairness and Machine Learning} (online book), selected sections.;%
Knaflic, \textit{Storytelling with Data}, selected chapters on communicating results.%
}
%
  }%
}


\title{Elements of AIML\\Unit 03 -- Lecture 02 Notes\\Train/Test/Validation Split \& Cross-Validation}
\author{Tofik Ali}
\date{\today}

\begin{document}
\maketitle
\tableofcontents

\section{Lecture Overview}
This lecture focuses on \textbf{evaluation methodology}: how to split data, how to tune models, and how to avoid common mistakes that produce misleading performance numbers.
\notesources{\AIMLUnitIIISources}

\section{How to use these notes (self-study)}
Whenever you see an evaluation number (accuracy/F1/RMSE), ask: ``Which data was this computed on, and was any information from test leaking into training?''.
This lecture gives you a repeatable workflow for honest evaluation that you should follow in labs and projects.

\section{Learning Outcomes}
\begin{itemize}[leftmargin=*]
  \item Explain train/validation/test splits and what decisions each supports.
  \item Describe k-fold cross-validation and stratified splitting.
  \item Identify data leakage and explain why it is dangerous.
  \item Use a basic scikit-learn split pattern correctly.
\end{itemize}

\section{Keywords}
\begin{itemize}[leftmargin=*]
  \item \textbf{Overfitting:} model fits noise in training data; poor generalization.
  \item \textbf{Train set:} used to fit model parameters.
  \item \textbf{Validation set:} used to tune hyperparameters / select model.
  \item \textbf{Test set:} used for final unbiased evaluation.
  \item \textbf{Cross-validation:} repeated train/validate splits to estimate performance.
  \item \textbf{Leakage:} information from outside training-time availability enters the model.
\end{itemize}

\section{Abbreviations and symbols}
\begin{itemize}[leftmargin=*]
  \item \textbf{CV}: cross-validation; $k$: number of folds in k-fold CV.
  \item \textbf{Train/Val/Test}: training, validation, and test datasets.
  \item \textbf{Metric}: a number summarizing performance (accuracy, F1, RMSE, etc.).
  \item $\hat{y}$: model predictions; $y$: true labels/targets.
  \item \textbf{Hyperparameter}: setting chosen by us (e.g., depth, $C$, $\lambda$), not directly learned as weights.
\end{itemize}

\section{Why splits are necessary}
If we evaluate on training data, we measure how well the model memorized the data, not how well it generalizes.
Splits simulate the ``future unseen'' setting.

\subsection{Generalization gap (quick intuition)}
Often training error is lower than test error:
\[
  \text{Train error} < \text{Test error}
\]
The difference is the \textbf{generalization gap}. Good evaluation tries to estimate test performance \emph{before} we deploy the model.

\section{Train/validation/test}
\subsection{Train set}
Used to learn parameters (e.g., weights in regression). We can run the training algorithm many times on the train set.

\subsection{Parameters vs hyperparameters}
\begin{itemize}[leftmargin=*]
  \item \textbf{Parameters} are learned from training data (e.g., weights $w$ in linear regression).
  \item \textbf{Hyperparameters} are chosen by us and control learning (e.g., number of neighbors $k$ in kNN, regularization $\lambda$, tree depth).
\end{itemize}
We typically tune hyperparameters using validation (or CV), not using the test set.

\subsection{Validation set}
Used for \textbf{model selection} and \textbf{hyperparameter tuning}. Examples of tuning decisions:
\begin{itemize}[leftmargin=*]
  \item choose model type (tree vs SVM),
  \item choose hyperparameters (regularization strength, depth),
  \item choose preprocessing options.
\end{itemize}

\subsection{Test set}
The test set is used \textbf{once} for the final report. If you repeatedly look at test results while tuning, the test set becomes part of training decisions and the estimate becomes optimistic.

\subsection{Typical split ratios (rule of thumb)}
There is no single best ratio, but common choices are:
\begin{itemize}[leftmargin=*]
  \item 80/20 (train/test) for quick experiments,
  \item 60/20/20 or 70/15/15 (train/val/test) when you need an explicit validation set.
\end{itemize}
When data is small, k-fold CV is often better than holding out a large validation set.

\section{k-fold cross-validation}
Cross-validation uses multiple train/validation splits to reduce variance in the estimate.
In k-fold CV:
\begin{itemize}[leftmargin=*]
  \item split data into $k$ folds,
  \item train on $k-1$ folds and validate on the remaining fold,
  \item repeat $k$ times and average.
\end{itemize}

\subsection{What CV is used for (two different goals)}
\begin{itemize}[leftmargin=*]
  \item \textbf{Model selection}: choose the best hyperparameters/model family.
  \item \textbf{Performance estimation}: estimate how well the chosen model will generalize.
\end{itemize}
In practice, we often do model selection with CV and then train one final model on the full training set.

\subsection{Stratified k-fold (classification)}
Stratification preserves class proportions across folds to avoid folds with too few positive examples.

\subsection{Other splitting variants (when needed)}
\begin{itemize}[leftmargin=*]
  \item \textbf{Group split / GroupKFold}: keep samples from the same person/customer in one split (avoids ``same-user'' leakage).
  \item \textbf{Time-series split}: training uses past, validation uses future (no shuffling).
  \item \textbf{Nested CV (concept)}: outer CV estimates generalization, inner CV tunes hyperparameters (reduces optimistic bias).
\end{itemize}

\section{Data leakage (examples)}
Leakage can happen in subtle ways:
\begin{itemize}[leftmargin=*]
  \item \textbf{Post-event features:} a column that is only known after the outcome.
  \item \textbf{Preprocessing leakage:} scaling/encoding computed on the full dataset.
  \item \textbf{Duplicate leakage:} same or near-identical examples in train and test.
\end{itemize}

\subsection{More leakage patterns to watch for}
\begin{itemize}[leftmargin=*]
  \item \textbf{Target leakage}: a feature is a direct proxy for the label (e.g., ``total\_payment\_received'' to predict ``loan default'').
  \item \textbf{Feature selection leakage}: selecting features using the whole dataset before CV (the validation folds influenced feature choice).
  \item \textbf{Preprocessing order mistake}: ``fit scaler'' on all data, then split.
\end{itemize}

\subsection{Correct preprocessing rule}
Fit preprocessing only on training data:
\begin{itemize}[leftmargin=*]
  \item compute mean/std for scaling on train,
  \item apply the same transformation to validation/test.
\end{itemize}

\section{Minimal code patterns}
\subsection{Train-test split}
\begin{lstlisting}[language=Python]
from sklearn.model_selection import train_test_split

X_train, X_test, y_train, y_test = train_test_split(
    X, y, test_size=0.2, random_state=42, stratify=y
)
\end{lstlisting}

\subsection{Stratified k-fold}
\begin{lstlisting}[language=Python]
from sklearn.model_selection import StratifiedKFold

cv = StratifiedKFold(n_splits=5, shuffle=True, random_state=42)
for train_idx, val_idx in cv.split(X, y):
    X_tr, X_val = X[train_idx], X[val_idx]
    y_tr, y_val = y[train_idx], y[val_idx]
\end{lstlisting}

\subsection{Leakage-safe pipelines (recommended)}
Use a pipeline so that preprocessing is learned only from training folds:
\begin{lstlisting}[language=Python]
from sklearn.pipeline import Pipeline
from sklearn.preprocessing import StandardScaler
from sklearn.linear_model import LogisticRegression

pipe = Pipeline([
    ("scaler", StandardScaler()),
    ("clf", LogisticRegression())
])
\end{lstlisting}

\section{Common pitfalls (and how to avoid them)}
\begin{itemize}[leftmargin=*]
  \item \textbf{Tuning on the test set:} if you look at test repeatedly, it stops being unbiased.
  \item \textbf{Preprocessing before splitting:} scaling/encoding on full data leaks information into test folds.
  \item \textbf{Wrong split strategy:} time-series and grouped data need special splits (no random shuffle).
  \item \textbf{Feature selection outside CV:} choosing features using all data leaks information into validation.
  \item \textbf{Reporting only one run:} randomness can change results; use CV or repeat splits when appropriate.
\end{itemize}

\section{Practice (with answers)}
\subsection*{P1. If you tune hyperparameters using the test set, what happens?}
\textbf{Answer:} the test set is no longer unbiased; you indirectly overfit to it, so the reported test score becomes overly optimistic.

\subsection*{P2. Give one example where time-series splitting is required.}
\textbf{Answer:} predicting next week's sales from past sales. Shuffling would mix future information into training.

\subsection*{P3. Why is GroupKFold useful?}
\textbf{Answer:} it prevents samples from the same group/person appearing in both train and validation/test, avoiding identity-based leakage.

\section{Exit questions (with answers)}

\subsection*{Q1. Why should the test set be used only at the end?}
\textbf{Answer:} because otherwise tuning decisions indirectly depend on test performance, causing optimistic bias; the test set is no longer an unbiased estimate.

\subsection*{Q2. What is cross-validation and when is it useful?}
\textbf{Answer:} CV repeats training/validation on multiple folds to get a more stable estimate; useful when data is limited or we want robust model selection.

\subsection*{Q3. Give two leakage examples.}
\textbf{Answer:} using a future field to predict the outcome; scaling using the full dataset before splitting; duplicates across train/test.

\subsection*{Q4. What is stratification?}
\textbf{Answer:} ensuring each split/fold has approximately the same class distribution as the full dataset, important for classification metrics.

\section{Summary}
Correct evaluation requires clean separation between training and final testing. Cross-validation and careful leakage prevention are core skills for reliable ML.
\notesources{\AIMLUnitIIISources}

\end{document}
